\documentclass[10pt]{article} % For LaTeX2e
\usepackage{tmlr}
% If accepted, instead use the following line for the camera-ready submission:
%\usepackage[accepted]{tmlr}
% To de-anonymize and remove mentions to TMLR (for example for posting to
% preprint servers), instead use the following:
%\usepackage[preprint]{tmlr}

\usepackage[utf8]{inputenc}
\usepackage{amsmath,amssymb,amsthm,mathtools}
\usepackage{booktabs}
\usepackage{float}
\usepackage{microtype}
\usepackage{pgfplots}
\usepackage{xcolor}
\usepackage[pdftitle={What a causal mask is worth: a rank-logit law for prefix averaging},
  colorlinks=true,linkcolor=blue!45!black,citecolor=blue!45!black,
  urlcolor=blue!45!black]{hyperref}
\usepackage{url}

\pgfplotsset{compat=1.18}
\newtheorem{theorem}{Theorem}
\newtheorem{proposition}[theorem]{Proposition}
\newtheorem{corollary}[theorem]{Corollary}
\newtheorem{lemma}[theorem]{Lemma}
\theoremstyle{definition}
\newtheorem{definition}[theorem]{Definition}
\DeclareMathOperator{\rank}{rank}
\DeclareMathOperator{\osc}{osc}
\newcommand{\R}{\mathbb R}
\newcommand{\one}{\mathbf 1}
\newcommand{\softmax}{\operatorname{softmax}}
\newcommand{\Stirling}[2]{\genfrac{\{}{\}}{0pt}{}{#1}{#2}}
\title{What a causal mask is worth:\\
a rank--logit law for prefix averaging\thanks{We used AI tools for proofreading
and improving the clarity of the manuscript.}}

% Authors must not appear in the submitted version. They are hidden by the
% tmlr package as long as neither [accepted] nor [preprint] is used.
\author{}

\def\month{08}  % Insert correct month for camera-ready version
\def\year{2026} % Insert correct year for camera-ready version
\def\openreview{\url{https://openreview.net/forum?id=XXXX}} % Insert correct link to OpenReview for camera-ready version

\begin{document}
\maketitle

\begin{abstract}
Transformer models use causal masks to force each position to assign zero
attention weight to future tokens. We ask what it costs for an unmasked dense
attention mechanism to approximate this behavior by giving sufficiently little
weight to those future tokens. We study this question in prefix averaging with
a single, position-dependent unmasked head. To achieve polynomial accuracy,
the logits required become extremely large. The smallest centered logit radius
\(B\) is governed by the law below, where \(p>1\) and positive integer \(r\)
are fixed, score matrices have centered score rank at most \(r\) and worst-row
output error is at most \(n^{-p}\):
\[
 \log B
 =\frac{n-1}{r}
 \left((p-1)\log n+\log\log n+O_{p,r}(1)\right).
\]

Matching lower and upper bounds show that this formula gives the optimal
asymptotic cost at every fixed polynomial error target \(n^{-p}\), including
both diverging terms. The lower bound uses a determinant argument, while the
upper bound constructs a balanced rank-\(r\) matrix. We also compare the
construction with these bounds numerically at finite sequence lengths. A
separate estimate shows that when \(p>3/2\), restricting the centered logit
radius to polynomial size forces the rank to grow proportionally with sequence
length.

These large logits force the product of the maximum query norm and maximum
centered-key norm to be large. For example, there are no query and key
coordinates inside float32's finite range that would achieve error \(n^{-2}\)
with a position-indexed pure dot-product head with dimension \(128\) once
\(n\ge4166\). The same fixed-rank law extends to a content-adaptive model in
which the Boolean values being averaged also affect the queries and keys.
\end{abstract}

\section{Introduction}

Attention is often described as a learned mechanism, where queries and keys
allow tokens to attend and influence one another. However, architectural
mechanisms also determine which key positions each query is allowed to attend
to before these scores are passed through softmax. In autoregressive
transformers, a causal mask supplies one such rule by preventing every
position from attending to future tokens. Causal masking is standard in
autoregressive transformers \citep{vaswani2017attention}. To our knowledge,
the literature has not yet quantified the cost of reproducing the same
operator without the mask.

We want to understand how much representational work this causal mask
performs. We do this by removing the causal mask and asking a dense attention
head to reproduce the same behavior using only its score matrix. Measuring the
additional score rank and logit magnitude required by an unmasked head allows
us to quantify the contribution of the causal mask, and also separates
structure supplied directly by architecture from structure that is represented
through learned parameters.

We study this question in the setting of prefix averaging, a standard
primitive in theoretical transformer constructions
\citep{yang2025cookbook,strobl2025}. In prefix averaging, a causal mask creates
a moving support boundary where all tokens after the current position are
masked and receive zero weight, while all visible tokens receive equal weight.
Let
\[
 T(i,j)=\frac1i\mathbf 1[j\le i]
\]
be the prefix-average operator. With causal masking, we can achieve \(T\)
exactly. With the mask removed, softmax forces us to assign positive mass to
every key. To approximate \(T\) with unmasked softmax, we'll require nearly
equal scores on each prefix along with strong suppression of every future key.

For an unmasked score matrix \(S\), define
\(P_0=I_n-\one\one^\top/n\), so that \(SP_0\) subtracts the average from each
row of \(S\). The centered score rank is \(\rho(S)=\rank(SP_0)\), and the
centered logit radius is \(B(S)=\max_{i,j}|(SP_0)(i,j)|\). For a pure
dot-product head of dimension \(d_h\), \(\rho(S)\le d_h\). Because softmax is
unchanged by row shifts, \(\softmax(S)=\softmax(SP_0)\). If \(\rho(S)\le r\),
then \(SP_0\) admits a factorization \(QK^\top\) using \(r\)-dimensional query
and key vectors. If two finite logit rows produce the same softmax
probabilities, they can only differ up to rowwise constants. Therefore,
\(\rho(S)\) is the smallest dimension needed to factorize any score matrix
that produces the same attention pattern. For every fixed \(p>1\) and fixed
positive integer \(r\), if \(B\) denotes the minimum centered logit radius
required to achieve worst-row output error \(n^{-p}\), our main theorem gives
\[
 r\log B
 =(n-1)\left((p-1)\log n+\log\log n+O_{p,r}(1)\right)
\]
as the optimum for error \(n^{-p}\).

This result says that a fixed-rank unmasked head must use centered logits whose
magnitude increases extremely quickly with sequence length. In this problem,
rank and logit magnitude can substitute for one another: extreme scores allow
a fixed-rank softmax to approximate the boundary, while higher rank supplies
more independent score patterns with which to represent the changing
boundary. A second result states that if the centered logit radius is
restricted to polynomial size and \(p>3/2\), then the rank must grow
proportionally with sequence length.

To achieve our main result, we use a determinant argument to show that no
rank-\(r\) score matrix can achieve the target with substantially smaller
logits. For the matching upper bound, we construct a balanced rank-\(r\)
matrix that achieves the same asymptotic scale.

We checked our result computationally. At selected values of \(n\) and \(r\),
we compare the realized centered logit radius of our construction with the
determinant lower bound and the construction's closed-form upper
budget. The construction lies between these bounds in every tested case. A
factor-norm inequality then converts the logit requirement into a lower bound
on the product of the maximum query norm and maximum centered-key norm. For
example, no position-indexed pure dot-product head of dimension \(128\) whose
query and key coordinates lie within
float32's finite range can achieve error \(n^{-2}\) once \(n\ge4166\). Under
the same assumptions, head dimensions through \(252\) are also insufficient
once \(n\ge8192\).

The results above hold one score matrix fixed while the value stream varies.
We next consider a restricted adaptive model in which the Boolean values being
averaged also affect the queries and keys. Uniform accuracy over all Boolean
inputs forces the attention matrix on the all-zero input to approximate causal
prefix attention, so the same fixed-rank law holds in this model.

To our knowledge, this is the first quantitative law for the rank--logit cost
of replacing a causal mask with dense unmasked attention in prefix averaging.
Our results suggest that structural encoding is a powerful
representational resource. When we know in advance that certain token pairs
shouldn't interact, explicitly masking those pairs may reduce the score rank
or logit magnitude that dense attention would otherwise require to approximate
zero attention. Our result also shows that rank alone can give an incomplete
picture of representational cost, as a low-rank head may be expressive enough
while requiring massive score magnitudes that are numerically unrepresentable
as sequence length grows.

The remainder of the paper is organized as follows. Section 2 discusses
related work. Section 3 sets up our problem with formal definitions for our
key mathematical objects. Section 4 presents our main rank--logit law and the
resulting rank--magnitude tradeoffs. Section 5 explains the determinant
mechanism behind the lower bound and the matching rank-\(r\) construction.
Section 6 presents a finite-size numerical audit that brackets the optimum and
derives bounds on query and key magnitudes. Section 7 extends the law to heads whose
queries and keys depend on the Boolean payload being averaged. Section 8
states three open questions. The appendices contain the complete proofs and
experimental details.

\section{Related work}

Transformer models combine learned mechanisms with architectural decisions.
In the original Transformer decoder, future positions are masked to preserve
the autoregressive property \citep{vaswani2017attention}. This represents a
different source of structure from learned dot-product scores.

Construction-oriented work uses causal masking to obtain prefix means or
prefix-sum primitives \citep{yang2025cookbook,strobl2025}; our result
quantifies the cost of replacing that masked primitive with one fixed-rank
dense softmax head.

Prior literature focuses on the behavior of the learned score mechanism.
\citet{bhojanapalli2020} show that head dimension can create a low-rank
bottleneck that prevents a head from representing some attention matrices.
\citet{amsel2025quality} show, for a nearest-neighbor target, that
adding many low-rank heads need not compensate for insufficient rank.
\citet{likhosherstov2023} show that sparse attention
matrices can be approximated with query/key dimension logarithmic in sequence
length when the error and sparsity parameters are fixed. These papers ask how constraints on rank or dimension
affect which attention patterns are representable. We instead look at one
attention pattern and study how large the logits need to become to approximate
it at fixed centered rank.

\citet{velickovic2025} show that when logit spread stays bounded, each
global softmax coefficient disperses on the order of \(1/n\) as the number of
items grows. Our result quantifies the logit magnitude
required to avoid such dispersion and relates it to score rank for causal
prefix averaging.

\citet{lee2026} studies row-centered logits as well under the name ``energy
field,'' and shows that tested trained models concentrate variance in
relatively few components. This suggests that a version of our result might
eventually be relevant to trained heads that behave approximately like
low-rank matrices. This remains a hypothesis.

\citet{alman2023} show that the time required to approximate dense attention
changes with bounds on the input entries. They ask how
quickly an attention output can be computed from given query and key factors;
we ask how large those factors must become to represent one target attention
pattern at fixed rank.

\section{Problem setup}

Let \(S\) denote the complete pre-softmax score matrix, including any
additive positional bias, and let \(A=\softmax_{\rm row}(S)\).  For a pure
dot-product head, we write
\[
 S(i,j)=\langle q_i,k_j\rangle,
\]
with the conventional \(1/\sqrt{d_h}\) scale absorbed into the factors.
All rank and radius bounds below apply to the complete matrix \(S\).
The relation \(\rho(S)\le d_h\) holds for pure dot-product scores; a
separately parameterized positional bias can add rank independently of
\(d_h\).
Softmax is invariant under row shifts, so
\(SP_0\) is the identifiable score matrix.  Its radius controls every
within-row score difference:
\[
 |S(i,j)-S(i,t)|\le 2B(S).
\]
For a bounded value stream \(v\), the output is \(Av\).  The operational
error is
\begin{equation}\label{eq:operator}
 \sup_{\|v\|_\infty\le1}\|(A-T)v\|_\infty
 =\max_i\|A(i,:)-T(i,:)\|_1 .
\end{equation}
H\"older's inequality gives the upper bound, and the sign vector of the
largest-error row gives equality.  The metric asks one realized attention
matrix to average every bounded value stream.

Uniformity ranges over separate value streams while \(A\) stays fixed.  A
single fixed value matrix tests the column span that it presents to the head.

The radius \(B(S)\) also admits an architectural bound.

\begin{proposition}[Factor-norm bound]\label{prop:norms}
Let \(S=QK^\top\) with rows \(q_i,k_j\in\R^{d_h}\), and let
\(\bar k=\frac1n\sum_jk_j\).  Then
\begin{equation}\label{eq:norms}
 B(S)\le\max_i\|q_i\|_2\cdot\max_j\|k_j-\bar k\|_2 .
\end{equation}
If in addition every entry of \(Q\) and \(K\) has magnitude at most \(F\),
then \(B(S)\le2d_hF^2\).  Consequently, if \(B(S)>0\) and \(F>0\),
then \(\log B(S)\le\log 2+\log d_h+2\log F\).
\end{proposition}

\begin{proof}
Centering a score row is equivalent to centering the keys: row \(i\) of
\(SP_0\) has entries
\(\langle q_i,k_j\rangle-\frac1n\sum_t\langle q_i,k_t\rangle
=\langle q_i,k_j-\bar k\rangle\), and Cauchy--Schwarz gives
\eqref{eq:norms}.  Under the entrywise bound,
\(\|q_i\|_2\le\sqrt{d_h}F\) and
\(\|k_j-\bar k\|_2\le2\sqrt{d_h}F\), hence
\(B(S)\le2d_hF^2\).  Taking logarithms gives the final claim when the
quantities are positive.
\end{proof}

Thus every lower bound on \(B(S)\) is also a lower bound on
\(\max_i\|q_i\|_2\max_j\|k_j-\bar k\|_2\).  If the conventional
\(1/\sqrt{d_h}\) scale is written explicitly rather than absorbed into the
factors, the right-hand side of \eqref{eq:norms} is divided by \(\sqrt{d_h}\).

Write \([m]=\{1,\ldots,m\}\).  For a finite set, \(\osc\) denotes the maximum
minus the minimum, and \(\Stirling{m}{k}\) denotes a Stirling number of the
second kind.  Let
\(I\subseteq[n-1]\) and put
\[
 I^+=I\cup\{j+1:j\in I\}.
\]
Row \(i\in I\) has a causal \((\delta_i,\gamma_i)\) profile
\emph{on \(I\)} when
\begin{equation}\label{eq:profile}
 \osc_{j\in I^+,\ j\le i}S(i,j)\le\delta_i,
 \qquad S(i,i+1)\le S(i,i)-\gamma_i.
\end{equation}
Accuracy sets explicit row-dependent parameters.

\begin{lemma}[Output error gives score geometry]\label{lem:profile}
Let \(i\in[n-1]\), let \(0<\varepsilon<1/(2i)\), and suppose
\(\max_j|A(i,j)-T(i,j)|\le\varepsilon\).  Then row \(i\) satisfies
\eqref{eq:profile} on every \(I\subseteq[n-1]\) containing \(i\), with
\begin{equation}\label{eq:rowprofile}
 \delta_i=\log\frac{1+i\varepsilon}{1-i\varepsilon},
 \qquad
 \gamma_i=\log\frac{1-i\varepsilon}{i\varepsilon}.
\end{equation}
\end{lemma}

This is a log-odds conversion: near-uniform prefix probabilities give
near-equal prefix scores, while the first future probability gives the
boundary gap.  Appendix~\ref{app:profile} gives the proof.

The lower bound uses only entrywise matrix error, which follows from the
operational row-\(\ell_1\) bound in \eqref{eq:operator}; the upper construction
attains the row-\(\ell_1\) target.  Therefore
\(\|x\|_\infty\le\|x\|_q\le\|x\|_1\) gives the same two-term law for maximum
row-\(\ell_q\) error at every fixed \(1\le q\le\infty\).

The scale \(1/n\) has a direct interpretation.  For a binary value stream,
\(i(Tv)_i\) is an integer prefix sum.  Rounding \(i(Av)_i\) recovers that sum
whenever the scalar output error is below \(1/(2i)\).  The eligibility
condition \(i\varepsilon<1/2\) marks this recovery scale.  Fixed targets
\(n^{-p}\), \(p>1\), give every nonterminal boundary (\(i\le n-1\)) a
polynomial margin.

\section{The rank--logit law}

Lemma~\ref{lem:profile} converts output accuracy into row profiles.  The
following nonasymptotic determinant bound converts those profiles into a
constraint on rank and logit radius, while Theorem~\ref{thm:upper} supplies
the matching construction.

\begin{theorem}[Row-subset determinant bound]\label{thm:subset}
Let \(\varnothing\ne I\subseteq[n-1]\), \(m=|I|\), and suppose the rows in
\(I\) satisfy \eqref{eq:profile} on \(I\), with positive
\(\delta_i,\gamma_i\).  If \(r\ge1\) is an integer,
\(\rho(S)\le r\) and \(B(S)\le B\), then
\begin{equation}\label{eq:weighted}
 \prod_{i\in I}\gamma_i
 \le
 \sum_{k=1}^{\min\{r,m\}}
 k!\Stirling{m}{k}(2B)^k
 \max_{\substack{R\subseteq I\\|R|=k}}
 \prod_{i\in I\setminus R}\delta_i .
\end{equation}
If additionally \(\delta_i\le\delta\) and
\(\gamma_i\ge\gamma\) on \(I\), then
\begin{equation}\label{eq:constant}
 \gamma^m\le
 \sum_{k=1}^{\min\{r,m\}}
 k!\Stirling{m}{k}\delta^{m-k}(2B)^k .
\end{equation}
\end{theorem}

Take adjacent score differences.  The selected matrix has low rank, a large
diagonal, and a small strict lower triangle.  The key move is to expand the
determinant of its upper-triangular part, rather than the determinant of the
low-rank matrix itself; Appendix~\ref{app:determinant} gives the proof.

The theorem can be applied to any subset of eligible rows, namely rows with
\(i\varepsilon<1/2\).  Our main result uses all \(n-1\) nonterminal rows,
which are eligible when \(\varepsilon=o(1/n)\).

For \(0<\varepsilon<1\) and an integer \(1\le r\le n-1\), define
\[
 B^{\rm out}_{r}(n,\varepsilon)=\inf\left\{B(S):
 \rho(S)\le r,\ 
 \sup_{\|v\|_\infty\le1}\|(A-T)v\|_\infty\le\varepsilon\right\},
 \qquad B^{\rm out}_{r,p}(n)=B^{\rm out}_{r}(n,n^{-p}).
\]
Write \(N=n-1\) through the rest of this section.

\begin{theorem}[Balanced rank-\(r\) construction]\label{thm:upper}
Let \(n\ge2\), \(0<\varepsilon\le1/n\), \(N=n-1\), and
let \(r\) be an integer with \(1\le r\le N\).  There is
\(S\in\R^{n\times n}\) with
\(\rank(S)=\rho(S)=r\), output error at most \(\varepsilon\), and
\begin{equation}\label{eq:upperbudget}
 \log B(S)\le
 \frac Nr\log\frac{64e\log(8n/\varepsilon)}{n\varepsilon}
 +\log\log\frac{8n}{\varepsilon}.
\end{equation}
\end{theorem}

The construction partitions the nonterminal rows into \(r\) balanced
consecutive blocks and uses one outer product per block.
Appendix~\ref{app:construction} verifies its rank, error, and radius.

\begin{theorem}[Two-term fixed-rank output law]\label{thm:sharp}
For every fixed \(p>1\) and fixed positive integer \(r\),
\begin{equation}\label{eq:sharp}
 \log B^{\rm out}_{r,p}(n)
 =\frac{n-1}{r}
 \bigl((p-1)\log n+\log\log n+O_{p,r}(1)\bigr).
\end{equation}
\end{theorem}

\begin{proof}
We first apply the determinant bound to every feasible score matrix and then
use Theorem~\ref{thm:upper} for the matching construction.

Set \(\varepsilon=n^{-p}\) and \(N=n-1\).  For the lower bound, fix any
feasible \(S\) in the definition of \(B^{\rm out}_{r,p}(n)\), put
\(B=B(S)\), and set
\[
 x_n=N\varepsilon,\qquad
 \delta_n=\log\frac{1+x_n}{1-x_n},\qquad
 \gamma_n=\log\frac{1-x_n}{x_n}.
\]
Since \(p>1\), \(x_n\to0\).  For all sufficiently large \(n\), operational
error at most \(\varepsilon\) and Lemma~\ref{lem:profile} give, uniformly on
\(I=[N]\),
\[
 \delta_i\le\delta_n=2n^{1-p}(1+o(1)),\qquad
 \gamma_i\ge\gamma_n=(p-1)\log n+o(1).
\]
Eventually \(\gamma_n\ge\delta_n\) and \(N\ge r\).  Since row centering
preserves score differences, \(S(i,i)-S(i,i+1)\ge\gamma_n\) in every
profiled row implies \(2B\ge\gamma_n\).  Hence, for \(k\le r\),
\[
 k!\Stirling{N}{k}\le k^N\le r^N,\qquad
 \delta_n^{N-k}(2B)^k
 \le\delta_n^{N-r}(2B)^r.
\]
Applying Theorem~\ref{thm:subset} and summing its at most \(r\) terms gives
\[
 \gamma_n^N\le r^{N+1}\delta_n^{N-r}(2B)^r,
\]
and therefore
\[
 \log B\ge\frac Nr\log\frac{\gamma_n}{r\delta_n}
 +\log\frac{\delta_n}{2}-\frac1r\log r.
\]
Here
\[
 \log\frac{\gamma_n}{r\delta_n}
 =(p-1)\log n+\log\log n+O_{p,r}(1),
\]
while, after normalization by \(N/r\), the remaining two terms are
\(O_{p,r}(\log n/n)=o(1)\).  Taking the infimum over all feasible \(S\)
proves the lower estimate; no minimizer is needed.

For the upper estimate, Theorem~\ref{thm:upper} gives
\[
 \log B\le\frac Nr
 \log\frac{64eL_n}{n\varepsilon}+\log L_n,
 \qquad L_n=\log\frac{8n}{\varepsilon}.
\]
At \(\varepsilon=n^{-p}\),
\[
 \log\frac{64eL_n}{n\varepsilon}
 =(p-1)\log n+\log\log n+O_p(1),
\]
and \((r/N)\log L_n=O_{p,r}(\log\log n/n)=o(1)\).  Combining the two
bounds proves \eqref{eq:sharp}.
\end{proof}

The \(O_{p,r}(1)\) term permits a multiplicative
\(\exp(O(n/r))\) uncertainty in \(B\).

\begin{corollary}[Head-dimension tradeoff]\label{cor:head}
Fix \(p>1\) and a positive integer \(d_h\).  A position-dependent unmasked
pure dot-product
head of dimension
\(d_h\) with operational output error at most \(n^{-p}\) obeys, for some constant
\(C_{p,d_h}\) and all sufficiently large \(n\),
\[
 d_h\log B(S)\ge
 (n-1)\bigl((p-1)\log n+\log\log n-C_{p,d_h}\bigr).
\]
\end{corollary}

\begin{proof}
Writing \(S=QK^\top\) gives
\(\rho(S)=\rank(SP_0)\le\rank(S)\le d_h\).  Thus \(S\) is feasible in the
definition of \(B^{\rm out}_{d_h,p}(n)\), and
\(B(S)\ge B^{\rm out}_{d_h,p}(n)\).  The lower bound in
Theorem~\ref{thm:sharp}, applied with the fixed rank parameter \(d_h\), gives
the displayed inequality.
\end{proof}

Theorem~\ref{thm:sharp} keeps \(r\) fixed.  If \(r\) grows, the upper bound
remains uniform for \(r=o(N/\log\log n)\), while the determinant lower bound
contains an additional \(-\log r\) term.  Thus the two diverging terms still
match when \(\log r=o(\log\log n)\).  Appendix~\ref{app:hadamard} gives a
complementary estimate for larger ranks.

At full centered rank, Proposition~\ref{prop:fullrank} attains
\(\log B=\log\log(n/\varepsilon)+O(1)\).  The closed-form bound
\eqref{eq:upperbudget} is loose at \(r=N\): its singleton-block construction
itself has \(\log B=\log\log n+O_p(1)\).

\begin{corollary}[Polynomial radius budget]\label{cor:budget}
Fix \(p>3/2\) and \(c>0\).  A score matrix with \(B(S)\le n^c\) and output
error at most \(n^{-p}\) has
\[
 \rho(S)\ge
 \left(\frac{p-3/2}{c+p-1}-o(1)\right)n.
\]
\end{corollary}

Appendix~\ref{app:hadamard} replaces the rank-dependent permutation count by
a columnwise Hadamard bound and proves the corollary uniformly over the rank.
The resulting estimate yields a positive coefficient
for \(p>3/2\).  Behavior at the threshold remains open.

The determinant bound separates the accuracy regimes more precisely than a
single asymptotic formula.

\begin{center}
\begin{tabular}{lll}
\toprule
accuracy \(\varepsilon\) & eligible rows & status for fixed \(r\) \\
\midrule
\(n^{-p}\), fixed \(p>1\) & all nonterminal
 & two-term law, Theorem~\ref{thm:sharp} \\
general \(o(1/n)\) & all nonterminal
 & two bounds with unmatched leading terms \\
\(c/n\), fixed \(c>0\) & linear subset
 & \(\Omega_{c,r}(n/r)\) to \(O_{c,r}(n\log\log n/r)\) \\
\(1/n\ll\varepsilon\ll1\) & \(\Theta(1/\varepsilon)\)
 & lower bound; construction open \\
\bottomrule
\end{tabular}
\end{center}

At \(\varepsilon=1/(n\log n)\), the determinant bound gives
\((N/r)(\log\log n+\log\log\log n+O(1))\).  The budget in
\eqref{eq:upperbudget} gives \((N/r)(2\log\log n+O(1))\).  Their leading
coefficients differ, which is why \eqref{eq:sharp} is restricted to
polynomial targets.  The behavior of \(B^{\rm out}_r\) at and above the
\(1/n\) scale remains open.

Full centered rank changes the scaling.

\begin{proposition}[Full-rank two-level scores]\label{prop:fullrank}
Let \(n\ge2\).  For \(\Gamma>0\), set
\[
 S_\Gamma(i,j)=-\Gamma\,\mathbf1[j>i].
\]
Then \(\rank(S_\Gamma P_0)=n-1\), \(B(S_\Gamma)\le\Gamma\), and every
nonterminal row \(i\) has a causal \((0,\Gamma)\) profile on every
\(I\subseteq[n-1]\) containing \(i\).  Taking
\(\Gamma=\log(2(n-1)/\varepsilon)\), for \(0<\varepsilon\le1\), gives
output error at most \(\varepsilon\).
\end{proposition}

\begin{proof}
The submatrix on rows \(1,\ldots,n-1\) and columns \(2,\ldots,n\) is
upper triangular with diagonal \(-\Gamma\), so \(\rank(S_\Gamma)\ge n-1\).
Column \(1\) of \(S_\Gamma\) is zero, since \(1>i\) fails for every
\(i\ge1\), so \(e_1\in\ker S_\Gamma\), the rank is \(n-1\), and
\(\ker S_\Gamma=\operatorname{span}\{e_1\}\).  If
\(S_\Gamma P_0x=0\), then \(P_0x\) lies in this kernel and has coordinate
sum zero, which gives \(P_0x=0\) and
\(\ker(S_\Gamma P_0)=\operatorname{span}\{\one\}\).  The centered rank is
\(n-1\).  Row centering keeps every entry within the raw row range
\(\Gamma\).  If \(\lambda_i\) is the future attention mass in row \(i\), then
\[
 \lambda_i=\frac{(n-i)e^{-\Gamma}}{i+(n-i)e^{-\Gamma}}
 \le(n-1)e^{-\Gamma}\le\varepsilon/2.
\]
Conditioned on the prefix, the attention distribution is uniform.  Its
row-\(\ell_1\) error therefore equals \(2\lambda_i\le\varepsilon\).
\end{proof}

The matrix \(S_\Gamma\) gives a finite-\(\Gamma\) approximation to the causal
mask.  It uses full centered rank and a growing realized radius to approximate
the architectural support boundary.

\section{Determinant mechanism and matching construction}\label{sec:mechanism}

Let \(N=n-1\), let \(\Delta\in\R^{n\times N}\) be the first-difference
operator, and take the first \(N\) rows of \(S\Delta\):
\[
 D(i,j)=S(i,j+1)-S(i,j).
\]
Since \(P_0\Delta=\Delta\), \(\rank(D)\le\rho(S)\).  Row \(i\)'s profile
gives a small strict lower triangle and a large diagonal:
\[
 |D(i,j)|\le\delta_i\quad(j<i),\qquad
 D(i,i)\le-\gamma_i.
\]
A direct determinant argument now fails: when \(\rank(D)<N\), it only says
\(\det D=0\).  Instead let \(U\) be the upper-triangular part of \(D\) and
write \(U=D-E\), where \(E\) is strictly lower triangular.  The determinant
of \(U\) is large because every diagonal entry \(U(i,i)=D(i,i)\) has
magnitude at least \(\gamma_i\).  In its column expansion, rank allows at
most \(r\) columns from \(D\); every other
column comes from the small lower triangle.  Lemma~\ref{lem:count} shows that
\(k!\Stirling{N}{k}\) pairs \((\pi,K)\) can contribute when \(k\)
columns are taken from \(D\).  Keeping the associated row indices yields
\eqref{eq:weighted}; Appendix~\ref{app:determinant}
supplies the complete argument, including arbitrary row subsets.

The construction balances row budgets over \(r\) consecutive blocks.  One
outer product creates the score drops inside each block.  Appendix~\ref{app:construction}
gives the factorization and verifies its centered rank,
row error, and radius bound \eqref{eq:upperbudget}.

\section{Numerical checks and query/key coordinate bounds}\label{sec:numerical}

The audit evaluates the explicit construction and the determinant bound at
\(p=2\).  High-precision decimal arithmetic computes the
row-\(\ell_1\) error and centered radius.  The lower curve solves
\eqref{eq:weighted} with the row profile from Lemma~\ref{lem:profile} and the
rank-independent bound.

Every quantity in the bound is measurable on a realized score
matrix: per-row oscillations and drops, the centered radius, and the centered
rank can be read off logits directly, so \eqref{eq:weighted} evaluates
at any \((n,r,\varepsilon)\) with no asymptotic step.  It gives a necessary,
not sufficient, condition for accuracy \(\varepsilon\) at the measured
resources.

For a radius value \(B\), define its normalized remainder at \(p=2\), the
value used throughout this section, by
\[
 R_{n,r}(B)=\frac{r}{n-1}\log B-\bigl(\log n+\log\log n\bigr),
\]
the general leading term being \((p-1)\log n+\log\log n\).
Theorem~\ref{thm:sharp} predicts \(R_{n,r}=O_r(1)\) for fixed \(r\).
Figure~\ref{fig:bracket} places the construction between the lower bound
and its closed-form budget for all 36 tested parameter pairs.  Every measured
row error is below \(0.251\,n^{-2}\).  The displayed interval bounds the
unknown optimum.

\begin{figure}[H]
\centering
\begin{tikzpicture}
\begin{axis}[
 width=.47\linewidth,height=4.4cm,
 xlabel={length \(n\)},ylabel={normalized remainder \(R_{n,r}\)},
 xmin=10,xmax=66,ymin=-4.2,ymax=9.6,ytick={-4,-2,0,2,4,6,8},
 grid=major,tick label style={font=\footnotesize},label style={font=\small},
 legend style={font=\scriptsize,at={(.5,1.02)},anchor=south,
               legend columns=3}]
\addplot+[mark=square*,black]
 table[x=n,y=upper_budget_remainder,col sep=comma,
 restrict expr to domain={\thisrow{rank}}{1:1}]
 {experiments/output/one_head_bracket.csv};
\addlegendentry{budget \(r{=}1\)}
\addplot+[mark=*,blue!70!black]
 table[x=n,y=constructed_remainder,col sep=comma,
 restrict expr to domain={\thisrow{rank}}{1:1}]
 {experiments/output/one_head_bracket.csv};
\addlegendentry{construction \(r{=}1\)}
\addplot+[mark=triangle*,orange!85!black]
 table[x=n,y=lower_remainder,col sep=comma,
 restrict expr to domain={\thisrow{rank}}{1:1}]
 {experiments/output/one_head_bracket.csv};
\addlegendentry{lower \(r{=}1\)}
\addplot+[mark=square,black,densely dotted]
 table[x=n,y=upper_budget_remainder,col sep=comma,
 restrict expr to domain={\thisrow{rank}}{8:8}]
 {experiments/output/one_head_bracket.csv};
\addlegendentry{budget \(r{=}8\)}
\addplot+[mark=o,blue!70!black,densely dotted]
 table[x=n,y=constructed_remainder,col sep=comma,
 restrict expr to domain={\thisrow{rank}}{8:8}]
 {experiments/output/one_head_bracket.csv};
\addlegendentry{construction \(r{=}8\)}
\addplot+[mark=triangle,orange!85!black,densely dotted]
 table[x=n,y=lower_remainder,col sep=comma,
 restrict expr to domain={\thisrow{rank}}{8:8}]
 {experiments/output/one_head_bracket.csv};
\addlegendentry{lower \(r{=}8\)}
\end{axis}
\end{tikzpicture}\hfill
\begin{tikzpicture}
\begin{axis}[
 width=.47\linewidth,height=4.4cm,
 xlabel={centered rank \(r\)},ylabel={\(\log B\) at \(n=64\)},
 xmode=log,log basis x=2,xtick={1,2,4,8},xmin=.85,xmax=9.5,ymin=0,ymax=800,
 grid=major,tick label style={font=\footnotesize},label style={font=\small},
 legend style={font=\scriptsize,at={(.98,.98)},anchor=north east}]
\addplot+[mark=triangle*,orange!85!black]
 table[x=rank,y=log_B_lower_weighted,col sep=comma,
 restrict expr to domain={\thisrow{n}}{64:64}]
 {experiments/output/one_head_bracket.csv};
\addlegendentry{determinant lower bound}
\addplot+[mark=*,blue!70!black]
 table[x=rank,y=log_B_constructed,col sep=comma,
 restrict expr to domain={\thisrow{n}}{64:64}]
 {experiments/output/one_head_bracket.csv};
\addlegendentry{radius of the constructed matrix}
\addplot+[mark=square*,black,dashed]
 table[x=rank,y=log_B_upper_budget,col sep=comma,
 restrict expr to domain={\thisrow{n}}{64:64}]
 {experiments/output/one_head_bracket.csv};
\addlegendentry{closed-form budget}
\end{axis}
\end{tikzpicture}
\caption{Finite-size construction and bracket at \(p=2\).  Left: the three
normalized remainders as \(n\) varies, with \(r=1\) solid and \(r=8\) dotted.
At \(n=64\), the full bracket spans \(6.12\) nats for \(r=1\) and \(9.25\)
for \(r=8\).  Right: the row-weighted determinant lower bound, radius of the
constructed matrix, and closed-form budget at \(n=64\).}
\label{fig:bracket}
\end{figure}

\subsection{Consequences of bounded query/key coordinates}

Proposition~\ref{prop:norms} converts the determinant lower bound into a statement
about heads whose query and key coordinates are bounded.

\begin{corollary}[Infeasibility under bounded query/key coordinates]\label{cor:precision}
Fix \(p>1\), a positive integer \(d_h\), and a floating-point format whose
largest finite magnitude is \(F>0\).  Let \(Q,K\in\R^{n\times d_h}\) have
entries bounded in magnitude by \(F\), and set \(S=QK^\top\).  Suppose the
operational output error is at
most \(\varepsilon=n^{-p}\).  For every nonempty
\(I\subseteq[n-1]\) on which \(i\varepsilon<1/2\), set \(m=|I|\) and
\[
 \delta_i=\log\frac{1+i\varepsilon}{1-i\varepsilon},\qquad
 \gamma_i=\log\frac{1-i\varepsilon}{i\varepsilon}.
\]
Then necessarily
\[
 \prod_{i\in I}\gamma_i
 \le\sum_{k=1}^{\min\{d_h,m\}}
 k!\Stirling{m}{k}(4d_hF^2)^k
 \max_{\substack{R\subseteq I\\|R|=k}}
 \prod_{i\in I\setminus R}\delta_i.
\]
\end{corollary}

\begin{proof}
Operational error implies entrywise error at most \(\varepsilon\), so
Lemma~\ref{lem:profile} supplies the displayed profiles.  Moreover
\(\rho(S)\le\rank(S)\le d_h\), while Proposition~\ref{prop:norms} gives
\(B(S)\le2d_hF^2\).  Substitution in Theorem~\ref{thm:subset} proves the
claim.
\end{proof}

Here the floating-point format is used only to set the coordinate bound \(F\);
the matrix product and softmax are evaluated in exact arithmetic.
Evaluation of the test at \(p=2\) gives the table.  Each entry is
the largest \(n\) not excluded by the bound.  The bfloat16 and float32
thresholds are computed from their respective maxima (\(\log F=88.7189\) and
\(88.7228\)); both round to the displayed value, and the nonexcluded lengths
coincide.

\begin{center}
\begin{tabular}{lccccc}
\toprule
format & \(\log F\) & \(d_h=8\) & \(32\) & \(64\) & \(128\) \\
\midrule
float16          & \(11.09\) & \(65\)  & \(262\)  & \(531\)  & \(1075\) \\
bfloat16         & \(88.72\) & \(268\) & \(1053\) & \(2093\) & \(4165\) \\
float32          & \(88.72\) & \(268\) & \(1053\) & \(2093\) & \(4165\) \\
\bottomrule
\end{tabular}
\end{center}

The first failing length remains excluded as \(n\) grows because each
\(\delta_i\) decreases and each \(\gamma_i\) increases on the same eligible
row subset.  At \(n=8192\), the bound gives \(\log B\ge418.72\)
for \(r=128\).  Proposition~\ref{prop:norms} then forces
\[
 \max_i\|q_i\|_2\max_j\|k_j-\bar k\|_2\ge e^{418.72}.
\]
The same calculation excludes head dimensions through \(252\) for
float32 or bfloat16 and through \(934\) for float16.

For scale, deployed decoder contexts run from \(1.3\times10^{5}\) to beyond
\(10^{6}\) positions \citep{llama3herd,gemini15}; every nonexcluded length in
the table lies below \(4.2\times10^{3}\), and the reference length \(n=8192\)
is a sixteenth of the shorter figure.

\section{Boolean query/key adaptation}\label{sec:adaptive}

The bounds so far hold one score matrix fixed while value streams vary.  In a
transformer, changing the input also changes the queries and keys, and hence
the attention matrix.  We ask whether changing this input can avoid the costs
imposed by our main theorem.  The answer is no in the following Boolean model.

Position \(j\) carries a payload bit \(v_j\in\{0,1\}\) and the embedding
\(x_j=e_j+v_ju_j\), with free position-indexed vectors \(e_j,u_j\in\R^d\).
A head computes \(q_i=Qx_i\), \(k_j=Kx_j\) with \(Q,K\in\R^{d_h\times d}\),
applies the unmasked row softmax \(A(v)\) of
\(S(v)_{ij}=\langle q_i,k_j\rangle\), and outputs \((A(v)v)_i\); the value
stream is the payload itself.  Each score \(S(v)_{ij}\) depends only on
\(v_i\) and \(v_j\).  In particular, when \(v_i=0\), changing \(v_j\) for
\(j\ne i\) changes only column \(j\) of row \(i\).
The head is \emph{\(\varepsilon\)-accurate} when
\(|(A(v)v)_i-\frac1i\sum_{j\le i}v_j|\le\varepsilon\) for every payload
\(v\in\{0,1\}^n\) and every nonterminal row \(i\).  Write
\(B_\infty=\max_v B(S(v))\); in particular \(B(S(0))\le B_\infty\).
Every realized score matrix factors through the head, so
\(\rho(S(v))\le d_h\) for every \(v\).

\begin{theorem}[Adaptive extraction]\label{thm:adaptive}
There are absolute constants \(i_0\in\mathbb N\) and \(c_0,C>0\) with the
following property.  If a head is \(\varepsilon\)-accurate, then every row
\(i_0\le i\le n-1\) with \(i\varepsilon\le c_0\) satisfies
\[
 \max_j|A(0)_{ij}-T(i,j)|\le C\varepsilon .
\]
\end{theorem}

The difficulty is that activating \(v_j\) changes both its value and its key,
so a Boolean input does not directly reveal \(A(0)_{ij}\).  We restrict to \(v_i=0\),
which fixes the query in row \(i\).  At each of two different prefix counts,
we compare Boolean inputs having the same number of active prefix bits.  Together,
these comparisons separate the contributions of the baseline and activated
scores.  Swapping one
active prefix bit shows that the baseline prefix weights are nearly equal;
summing over cyclic windows determines their total; and activating all future
bits bounds the future mass.
Lemma~\ref{lem:profile} then converts the resulting entrywise estimate into
the score profile required by Theorem~\ref{thm:subset}.
Appendix~\ref{app:adaptive} gives the proof with explicit constants.

For \(0<\varepsilon<1\) and a positive integer \(r\), let
\(B^{\rm ad}_{r}(n,\varepsilon)\) denote the infimum of \(B_\infty\) over
all \(\varepsilon\)-accurate heads with \(d_h\le r\), allowing \(d\) and
all \(e_j,u_j,Q,K\) to vary.  Set
\(B^{\rm ad}_{r,p}(n)=B^{\rm ad}_r(n,n^{-p})\).

\begin{corollary}[Fixed-rank law under Boolean adaptation]\label{cor:adaptive}
For every fixed \(p>1\) and fixed positive integer \(r\),
\[
 \log B^{\rm ad}_{r,p}(n)
 =\frac{n-1}{r}\bigl((p-1)\log n+\log\log n+O_{p,r}(1)\bigr).
\]
\end{corollary}

The lower bound applies Theorem~\ref{thm:subset} to the all-zero score matrix
\(S(0)\), whose centered rank is at most \(r\); uniform accuracy over Boolean
inputs is used only to force the required profile on \(S(0)\).  The upper
bound embeds the fixed construction of Theorem~\ref{thm:upper} with \(u_j=0\).
For this model,
adaptation leaves both diverging terms of the fixed-rank law unchanged.
Appendix~\ref{app:adaptive} gives the proof.

\section{Open problems}\label{sec:scope}

The fixed-rank law leaves three main questions.  First, what happens when rank
grows with sequence length?  Resolving this would also determine whether
polynomially bounded logits force rank proportional to \(n\) throughout
\(p>1\), rather than only for \(p>3/2\).  Second, what is the correct logit
cost when the error is of order \(1/n\), where our lower and upper bounds do
not match?  Third, can rank be replaced by a noise-stable notion of
effective rank, so that the bound can be applied to score matrices measured
in trained models?

\bibliography{audited_references}
\bibliographystyle{tmlr}

\appendix

\section{Output error and score geometry}\label{app:profile}

\begin{proof}[Proof of Lemma~\ref{lem:profile}]
For \(j\le i\), \(T(i,j)=1/i\), so entrywise error at most \(\varepsilon\)
gives
\[
 \frac{1-i\varepsilon}{i}\le A(i,j)\le
 \frac{1+i\varepsilon}{i}.
\]
Score differences are logarithms of probability ratios.  In particular, for
\(j,t\le i\),
\[
 S(i,j)-S(i,t)=\log\frac{A(i,j)}{A(i,t)},
 \qquad
 |S(i,j)-S(i,t)|\le\log\frac{1+i\varepsilon}{1-i\varepsilon}.
\]
This bounds the oscillation on every subset of prefix columns by
\(\delta_i\).  For the adjacent key,
\(T(i,i+1)=0\) gives \(A(i,i+1)\le\varepsilon\), while
\(A(i,i)\ge(1-i\varepsilon)/i\), which yields
\[
 S(i,i+1)-S(i,i)
 =\log\frac{A(i,i+1)}{A(i,i)}
 \le-\log\frac{1-i\varepsilon}{i\varepsilon}.
\]
\end{proof}

\section{Determinant proof}\label{app:determinant}

\begin{lemma}[Permitted permutation count]\label{lem:count}
For \(1\le k\le m\),
\[
 \#\bigl\{(\pi,K):K\subseteq[m],\ |K|=k,\ \pi\in\mathfrak S_m,\
 \ \pi(j)>j\text{ for all }j\notin K\bigr\}
 =k!\Stirling{m}{k}.
\]
\end{lemma}

\begin{proof}
We construct a bijection between the counted pairs \((\pi,K)\) and pairs
consisting of a partition of \([m]\) into \(k\) blocks and a bijection from
the block maxima to the block minima.

The largest element \(j\) of any cycle of \(\pi\) has \(\pi(j)\le j\) and
hence lies in \(K\), so iterating \(\pi\) from any point meets \(K\).  Let
\(f(x)\) be the first element of \(K\) on the orbit
\(x,\pi(x),\pi^2(x),\ldots\); thus \(f(x)=x\) on \(K\).  The fibers of \(f\)
partition \([m]\) into \(k\) blocks, and the block \(F_a=f^{-1}(a)\) meets
\(K\) in \(\{a\}\).

Fix \(a\in K\) and write \(F_a=\{y_1<\cdots<y_t\}\).  Each
\(x\in F_a\setminus\{a\}\) lies outside \(K\), so \(\pi(x)>x\); and
\(\pi(x)\in F_a\), since either \(\pi(x)=a\) or \(\pi(x)\notin K\) and its
orbit first meets \(K\) at \(a\).  So \(\pi\) maps \(F_a\setminus\{a\}\)
injectively into \(F_a\) with \(\pi(x)>x\).  The largest element \(y_t\)
admits no such image, so \(a=y_t=\max F_a\).  The domain
\(F_a\setminus\{a\}\) and the set \(\{y_2,\ldots,y_t\}\) both have
\(t-1\) elements, and \(\pi(x)>x\) prevents \(y_1\) from being an image.
Thus \(\pi\) maps the former set bijectively onto the latter.  If \(t=1\),
there is nothing more to show.  If \(t\ge2\), then
\(\pi(y_{t-1})=y_t\).  Descending induction gives the entire chain: once
\(y_{j+2},\ldots,y_t\) are the images of \(y_{j+1},\ldots,y_{t-1}\), the
only unused image larger than \(y_j\) is \(y_{j+1}\).  Hence
\(\pi(y_j)=y_{j+1}\).
Thus off \(K\) the permutation is the successor map of the blocks, its
images are the nonminimal elements of the blocks, and \(\pi\) must
map \(K\), the set of block maxima, bijectively onto the set of block
minima.

Conversely, let \(\Pi\) be a partition of \([m]\) into \(k\) blocks and
\(\sigma\) a bijection from the block maxima onto the block minima.  Take
\(K\) to be the set of block maxima and let \(\pi\) act as the successor
map within blocks on \([m]\setminus K\) and as \(\sigma\) on \(K\).  Every
nonminimal element has one preimage, its block predecessor, and
every block minimum has one under \(\sigma\); so
\(\pi\in\mathfrak S_m\), and \(\pi(j)>j\) off \(K\).  From any nonmaximal
element the successor chain climbs within its block and meets \(K\) first
at the block maximum, so the fibers of \(f\) recover \(\Pi\), and the
restriction of \(\pi\) to \(K\) recovers \(\sigma\).  The two constructions
are therefore mutually inverse, and the pairs \((\pi,K)\) are in bijection
with the pairs \((\Pi,\sigma)\).  Since \(\Stirling{m}{k}\) counts the
partitions of \([m]\) into \(k\) nonempty blocks, there are
\(k!\Stirling{m}{k}\) such pairs.
\end{proof}

\begin{proof}[Proof of Theorem~\ref{thm:subset}]
Let \(N=n-1\), let \(\Delta\in\R^{n\times N}\) be the first-difference
operator, and set
\[
 D=(SP_0)\Delta=S\Delta,
 \qquad D(i,j)=S(i,j+1)-S(i,j).
\]
Each column of \(\Delta\) sums to zero, so \(P_0\Delta=\Delta\); the
representation \(D=(SP_0)\Delta\) gives the rank and entry bounds, while
\(D=S\Delta\) gives the formula for \(D(i,j)\).
Then \(\rank(D)\le\rho(S)\), and every entry of \(D\) has magnitude at most
\(2B\).  Order \(I=\{i_1<\cdots<i_m\}\) and define
\(D_I=(D(i_a,i_b))_{a,b=1}^m\).  This submatrix has rank at most \(r\).
For \(b<a\), both \(i_b\) and \(i_b+1\) lie in \(I^+\) and are at most
\(i_a\); recall that the profile bounds the oscillation of \(S(i,j)\) over
\(j\in I^+\), \(j\le i\), by \(\delta_i\) and gives
\(S(i,i+1)\le S(i,i)-\gamma_i\), so the profile of row \(i_a\) yields
\[
 |D_I(a,b)|\le\delta_{i_a},
 \qquad D_I(a,a)\le-\gamma_{i_a}.
\]

Let \(U_I\) be the upper-triangular part of \(D_I\), and let
\(E_I=D_I-U_I\).  Thus \(E_I\) is strictly lower triangular and
\[
 |\det U_I|\ge\prod_{i\in I}\gamma_i,\qquad U_I=D_I-E_I.
\]
Expand \(\det U_I=\det(D_I-E_I)\) multilinearly in its \(m\) columns.  For
\(K\subseteq[m]\), take column \(j\) from \(D_I\) when \(j\in K\) and from
\(E_I\) otherwise; write \(M_K\) for the resulting matrix.  Signs will be
discarded after taking absolute values.
If \(|K|>r\), the selected \(D_I\)-columns are linearly dependent, so the
mixed determinant vanishes.  If \(K=\varnothing\), the mixed matrix is
strictly lower triangular and its determinant also vanishes.

Fix \(|K|=k\).  In the Leibniz expansion, a product can be nonzero only when
\(\pi(j)>j\) for every \(j\notin K\), because column \(j\) of \(E_I\) is
supported strictly below row \(j\).  Lemma~\ref{lem:count} counts the pairs
\((\pi,K)\) with this property.

Every such product uses \(k\) entries from \(D_I\), each bounded by \(2B\).
An entry from \(E_I\) in row \(a\) is bounded by \(\delta_{i_a}\).  A
determinant product uses each row once, so if its \(D_I\)-entries occupy the
original rows in a \(k\)-element set \(R\subseteq I\), its
\(E_I\)-entries contribute at most
\[
 \max_{\substack{R\subseteq I\\|R|=k}}
 \prod_{i\in I\setminus R}\delta_i .
\]
Summing absolute values proves \eqref{eq:weighted}; replacing all
\(\delta_i,\gamma_i\) by their common bounds proves
\eqref{eq:constant}.
\end{proof}

\section{Hadamard bound and rank-independent radius lower bound}
\label{app:hadamard}

\begin{proposition}[Hadamard bound and rank-independent radius lower bound]\label{prop:hadamard}
Under the row-dependent hypotheses of Theorem~\ref{thm:subset},
\[
 B\ge\frac12\max_{i\in I}\gamma_i .
\]
In particular, under the constant-profile hypotheses, \(B\ge\gamma/2\).
If, in addition, \(m=|I|\), \(\delta\le\gamma\), and \(r\le m\), then
\[
 \gamma^m\le
 2^{m-1}(\delta\sqrt m)^{m-r}(2B\sqrt m)^r .
\]
\end{proposition}

\begin{proof}
For the radius bounds, row centering preserves score differences, so
every profiled row satisfies
\[
 \gamma_i\le |(SP_0)(i,i)-(SP_0)(i,i+1)|\le2B(S)\le2B .
\]
Maximizing over \(i\in I\) proves the general statement and its
constant-profile specialization.

For the determinant bound, assume \(\delta\le\gamma\) and \(r\le m\),
so that the radius bound gives \(2B\ge\gamma\ge\delta\).  Use the column
expansion from the proof of Theorem~\ref{thm:subset}.  As there,
\(\gamma^m\le|\det U_I|\le\sum_K|\det M_K|\), the sum over column
subsets \(K\) drawn from \(D_I\).  Column \(m\) of
\(E_I\) is zero, so every term taking column \(m\) from \(E_I\)
vanishes; terms with \(|K|>r\) vanish by rank.  At most \(2^{m-1}\)
terms remain, one for each set \(K\ni m\) with \(|K|\le r\).  In the
ordered \(m\)-by-\(m\) submatrix, every column of \(E_I\) has at most
\(m-1\) possibly nonzero entries, each bounded by \(\delta\), hence
norm at most \(\delta\sqrt m\); every column of \(D_I\) has norm at
most \(2B\sqrt m\).  For a term using the columns \(K\) from \(D_I\),
with \(|K|=k\), Hadamard's inequality gives
\[
 |\det M_K|
 \le(\delta\sqrt m)^{m-k}(2B\sqrt m)^{k}
 \le(\delta\sqrt m)^{m-r}(2B\sqrt m)^{r},
\]
the second step because \(k\le r\) and \(\delta\le2B\).  Since
\(|\det U_I|\ge\prod_{i\in I}\gamma_i\ge\gamma^m\), summing the
remaining terms proves the determinant bound.
\end{proof}

\begin{proof}[Proof of Corollary~\ref{cor:budget}]
Set \(N=n-1\), \(\varepsilon=n^{-p}\), \(r=\rho(S)\), and \(B=B(S)\).
Operational error implies entrywise error at most \(\varepsilon\).  In the
row-dependent profile supplied by Lemma~\ref{lem:profile}, \(\delta_i\)
increases and \(\gamma_i\) decreases with \(i\).  Thus, for all sufficiently
large \(n\), the values at \(i=N\) give the following common profile on
\(I=[N]\):
\[
 \delta=\log\frac{1+N\varepsilon}{1-N\varepsilon},\qquad
 \gamma=\log\frac{1-N\varepsilon}{N\varepsilon}.
\]
Here \(\gamma>0\) forces \(r\ge1\), and
\(r\le N\) holds unconditionally; eventually \(\gamma\ge\delta\).
Proposition~\ref{prop:hadamard} then gives \(2B\ge\gamma\ge\delta\) and
\[
 \gamma^N\le
 2^{N-1}(\delta\sqrt N)^{N-r}(2B\sqrt N)^r .
\]
Taking logarithms and isolating the rank-dependent term gives
\[
 r\log\frac{2B}{\delta}
 \ge N\log\frac{\gamma}{\delta\sqrt N}-(N-1)\log2.
\]
The profile parameters satisfy
\[
 \delta=2n^{1-p}(1+o(1)),\qquad
 \gamma=(p-1)\log n+O_p(1),\qquad
 \sqrt N=n^{1/2}(1+o(1)).
\]
Therefore
\[
 \log\frac{\gamma}{\delta\sqrt N}
 =(p-3/2)\log n+\log\log n+O_p(1).
\]
On the other hand, \(B\le n^c\) and
\(1/\delta=\tfrac12n^{p-1}(1+o(1))\) give
\[
 0\le\log\frac{2B}{\delta}
 \le(c+p-1)\log n+O_p(1).
\]
Replacing the nonnegative multiplier of \(r\) by this upper bound can only
weaken the rank lower bound.  Dividing the resulting inequality by
\(N\log n\), and using \(N/n\to1\), proves
\[
 r\ge\left(\frac{p-3/2}{c+p-1}-o(1)\right)n.
\]
All remainders depend only on the fixed \(p,c\), so this estimate is uniform
in the possibly \(n\)-dependent rank \(r\).
\end{proof}

\section{The upper construction}\label{app:construction}

\begin{proof}[Proof of Theorem~\ref{thm:upper}]
We choose a small allowable variation \(\delta_i\) among the prefix scores
and a gap \(\gamma_i\) separating the future scores.  One outer product will
handle each block of rows, and the partition will balance the largest radius
produced by these factors.

For \(1\le i\le N\), set
\[
 \delta_i=\frac{i\varepsilon}{16},\qquad
 \gamma_i=\log\frac{8(n-i)}{i\varepsilon},\qquad
 w_i=\log\left(1+\frac{\gamma_i}{\delta_i}\right),
\]
and write \(L_\varepsilon=\log(8n/\varepsilon)\).  Using
\(i\varepsilon\le N/n<1\le n-i\) and \(n-i\le in\), every row satisfies the
following bounds, which will be used repeatedly:
\begin{equation}\label{eq:margins}
 \delta_i\le\frac1{16}<\log8\le\gamma_i\le L_\varepsilon,
 \qquad
 (n-i)e^{-\gamma_i}=\frac{i\varepsilon}8 .
\end{equation}
Empty sums are zero.

\paragraph{Balanced partition.}
The row budget grows multiplicatively by
\(1+\gamma_t/\delta_t\), so equal-length blocks can have very different
radii.  We instead balance their logarithmic costs \(w_t\), which controls
the largest block product.
Partition \([N]\) into \(r\) consecutive
nonempty blocks so that, for every \(I=\{a,\ldots,b\}\),
\[
 \sum_{t=a+1}^{b}w_t\le W/r,\qquad W=\sum_{i=1}^Nw_i.
\]
A left-to-right greedy scan through \(w_2,\ldots,w_N\) starts a new block
before \(t\) whenever including \(w_t\) would make the current block's
internal load exceed \(W/r\).  At such a cut, the completed internal load
together with \(w_t\) exceeds \(W/r\).  The weight \(w_t\) becomes the first
index of the new block and is not counted in that block's internal load, so
the groups associated with different cuts are disjoint.  Their total weight
is at most \(W\), and therefore there are fewer than \(r\) cuts.
If the resulting partition has fewer than
\(r\) blocks, \(r\le N\) guarantees a nonsingleton block.  Splitting such
blocks reaches \(r\) and only removes internal weights.

\paragraph{Construction.}
For one block, \(v^I\) accumulates the required score drops and \(u^I\)
rescales each row.  Define
\[
 c_a=1,\qquad G_t=\sum_{s=a}^{t}\gamma_sc_s,\qquad
 c_t=\frac{G_{t-1}}{\delta_t}\quad(t>a),
\]
\[
 u_i^I=\frac{\mathbf1[i\in I]}{c_i},\qquad
 v_j^I=-\sum_{t=a}^{\min\{b,j-1\}}\gamma_tc_t,
\qquad
 S=\sum_Iu^I(v^I)^\top .
\]
The definitions form a joint recursion in increasing \(t\):
\(G_a=\gamma_a\), and \(G_t=G_{t-1}(1+\gamma_t/\delta_t)\) for \(t>a\).
By \eqref{eq:margins}, consecutive \(G\)'s grow by the factor
\(1+\gamma_t/\delta_t\ge2\), and \(c_i\ge1\) because
\(\delta_i\le1/16<\log8\le G_a\le G_{i-1}\) for \(i>a\).  For \(i\in I\),
row \(i\) of \(S\) is \(v^I/c_i\):
\[
 S(i,j)=
 \begin{cases}
  0,&j\le a,\\
  -G_{j-1}/c_i,&a<j\le b+1,\\
  -G_b/c_i,&j>b+1 .
 \end{cases}
\]
If \(i=a\), all prefix scores \(S(i,j)\), \(j\le i\), are zero.  If
\(i>a\), they decrease from \(0\) to
\(S(i,i)=-G_{i-1}/c_i=-\delta_i\).  In either case, every future score
\(S(i,j)\), \(j>i\), is at most
\(-G_i/c_i=S(i,i)-\gamma_i\le-\gamma_i\).  Row \(n\) of \(S\) is zero.

\paragraph{Rank.}
\(S\) is a sum of \(r\) outer products, so \(\rho(S)\le\rank(S)\le r\).
Suppose \(\sum_Ia_IP_0v^I=0\).  Then \(\sum_Ia_Iv^I\) lies in \(\ker P_0\),
so it is constant.  Fix a block \(I\) with first row \(a\).  Across
coordinates \(a\) and \(a+1\), every \(v^J\) with \(J\ne I\) is unchanged,
while \(v^I_{a+1}-v^I_a=-\gamma_ac_a\).  Constancy therefore gives
\[
 0=\sum_Ja_J(v^J_{a+1}-v^J_a)=-a_I\gamma_ac_a,
\]
and \(\gamma_ac_a>0\) forces \(a_I=0\).  Thus the vectors \(P_0v^I\) are
linearly independent.  At the first row \(a\) of block \(I\), row \(a\) of
\(S\) is \((v^I)^\top\), because \(c_a=1\).  Hence the corresponding row of
\(SP_0\) is \((P_0v^I)^\top\), so \(\rho(S)\ge r\).  Therefore
\(\rank(S)=\rho(S)=r\).

\paragraph{Output error.}
Fix \(i\in I\), let \(y_j=e^{S(i,j)}\le1\), and split the softmax mass of
row \(i\) into \(Z'=\sum_{j\le i}y_j\) and \(F=\sum_{j>i}y_j\), so
\(A(i,j)=y_j/Z\) with \(Z=Z'+F\).  The bound on the future scores and the
identity \((n-i)e^{-\gamma_i}=i\varepsilon/8\) give
\(F\le i\varepsilon/8\), while
\(1-e^{-x}\le x\) and the geometric growth of \(G\) bound the prefix
deficit.  If \(i=a\), then \(i-Z'=0\).  If \(i>a\), then
\[
 i-Z'=\sum_{j\le i}\bigl(1-e^{S(i,j)}\bigr)\le\sum_{j\le i}|S(i,j)|
 =\frac1{c_i}\sum_{t=a}^{i-1}G_t\le\frac{2G_{i-1}}{c_i}=2\delta_i
 =\frac{i\varepsilon}8,
\]
because each \(G_t\) is at most half of \(G_{t+1}\).  Thus
\(i-Z'\le i\varepsilon/8\) in every row.  Rows \(A(i,:)\) and \(T(i,:)\)
are probability vectors, so the error is twice the total undershoot.  Write
\(x_+=\max\{x,0\}\).  Since \(T(i,j)=0\) for \(j>i\), only prefix indices
contribute to the positive part:
\[
 \|A(i,:)-T(i,:)\|_1=2\sum_j\bigl(T(i,j)-A(i,j)\bigr)_+
 =2\sum_{j\le i}\Bigl(\frac1i-\frac{y_j}Z\Bigr)_+ .
\]
Put \(\bar Z=i(1+\varepsilon/8)\).  Since \(Z'\le i\) and
\(F\le i\varepsilon/8\), we have \(Z\le\bar Z\), so each undershoot is at
most \(1/i-y_j/\bar Z\), which is nonnegative because
\(y_j\le1\le\bar Z/i\).  Summing over the prefix,
\[
 \|A(i,:)-T(i,:)\|_1
 \le2\Bigl(1-\frac{Z'}{\bar Z}\Bigr)
 =\frac{2\bigl((i-Z')+i\varepsilon/8\bigr)}{\bar Z}
 \le\frac2i\cdot\frac{i\varepsilon}4=\frac\varepsilon2 .
\]
Row \(n\) is zero in \(S\), so its softmax is row \(n\) of \(T\),
and the worst-row error is at most \(\varepsilon/2\).

\paragraph{Radius bound.}
Unrolling the recursion across a block and using \(\gamma_a\le
L_\varepsilon\) with the balanced loads,
\[
 \log G_b=\log\gamma_a+\sum_{t=a+1}^bw_t\le\log L_\varepsilon+\frac Wr .
\]
Since \(\gamma_t/\delta_t\ge1\) by \eqref{eq:margins}, each
\(w_t\le\log(2\gamma_t/\delta_t)\le\log\frac{32L_\varepsilon}{t\varepsilon}\),
so \(N!\ge(N/e)^N\) and \(n\le2N\) give
\[
 W\le N\log\frac{32L_\varepsilon}\varepsilon-\log N!
 \le N\log\frac{32eL_\varepsilon}{N\varepsilon}
 \le N\log\frac{64eL_\varepsilon}{n\varepsilon} .
\]
Every \(c_i\ge1\), so all entries of a block row \(v^I_j/c_i\) lie in
\([-G_b,0]\); subtracting the row mean leaves every magnitude at most the
row range \(G_b\), and row \(n\) is zero.  Hence \(B(S)\le\max_IG_b\); the
block bound on \(\log G_b\) and the weight bound on \(W\) give
\eqref{eq:upperbudget}.
\end{proof}

\section{Adaptive extraction}\label{app:adaptive}

This appendix proves Theorem~\ref{thm:adaptive} and
Corollary~\ref{cor:adaptive}, and transfers the polynomial radius budget of
Corollary~\ref{cor:budget} to the adaptive model
(Corollary~\ref{cor:adaptivebudget}).  The proof gives
\(i_0=12\), \(c_0=10^{-3}\), and \(C=74\) in
Theorem~\ref{thm:adaptive}; these constants are not optimized.

\begin{proof}[Proof of Theorem~\ref{thm:adaptive}]
We first express every Boolean input using the scores from the all-zero input
and the scores obtained after activating individual positions.  Comparisons
at two prefix counts then show that the baseline prefix probabilities are
close to \(1/i\), while a final comparison bounds the probability assigned to
future positions.

Take \(i_0=12\) and \(c_0=10^{-3}\).  Fix a row \(i\) with
\(12\le i\le n-1\) and use only payloads
\(v=\one_X\) with \(X\subseteq[n]\setminus\{i\}\).  Every such payload has
\(v_i=0\), so the query \(q_i=Qe_i\) remains fixed, while key \(j\) takes
two possible values: the baseline \(Ke_j\) when \(v_j=0\) and the switched
\(K(e_j+u_j)\) when \(v_j=1\).  Thus row \(i\) has two possible scores in
each column, selected by \(v_j\).  Using the normalizing constant from the
all-zero input, set
\[
 Z=\sum_t e^{S(0)_{it}},\qquad
 p_j=\frac{e^{S(0)_{ij}}}{Z},\qquad
 \tilde p_j=\frac{e^{S(\one_{\{j\}})_{ij}}}{Z}\quad(j\ne i),
\]
put \(P_X=\sum_{j\in X}p_j\) and \(\widetilde P_X=\sum_{j\in X}\tilde p_j\).
Here \(p_j=A(0)_{ij}\), while \(\tilde p_j\) uses the score obtained when
\(v_j=1\) but retains the normalizing constant \(Z\) from the all-zero input.
Thus \(\tilde p_j\) need not be a probability.  Every \(p_j\) and
\(\tilde p_j\) is positive, and
\(1-P_X=\sum_{t\notin X}p_t\ge p_i>0\) for every admissible \(X\).
On the payload \(\one_X\), the exponentiated scores sum to
\(Z\widetilde P_X\) over \(X\) and to \(Z(1-P_X)\) outside \(X\).
Since the values equal one precisely on \(X\), the row-\(i\) output is
\begin{equation}\label{eq:face}
 F(X)=\frac{\widetilde P_X}{(1-P_X)+\widetilde P_X}.
\end{equation}
Accuracy states \(|F(X)-k/i|\le\varepsilon\) with
\(k=|X\cap[i-1]|\), since \(v_i=0\) removes bit \(i\) from the prefix
sum.  Only the prefix count of \(X\) enters the target; \(X\) may
contain future positions.

\paragraph{A linear constraint at fixed prefix count.}
Write \(Z_X=(1-P_X)+\widetilde P_X\) for the denominator of
\eqref{eq:face} and suppose \(k\le 2i/3\).  Multiplying the accuracy
bound by \(Z_X>0\) gives
\(\bigl|\widetilde P_X-\tfrac kiZ_X\bigr|\le\varepsilon Z_X\), and
expanding \(Z_X\) rearranges the left side:
\[
 \widetilde P_X-\tfrac ki\bigl((1-P_X)+\widetilde P_X\bigr)
 =\tfrac{i-k}i\Bigl(\widetilde P_X+h_kP_X-h_k\Bigr),
 \qquad h_k=\frac k{i-k}.
\]
It remains to bound \(Z_X\).  Since \(i\varepsilon\le c_0\le1\) and
\(i\ge12\) force \(\varepsilon\le\tfrac1{12}\), the output obeys
\(F(X)\le k/i+\varepsilon\le\tfrac23+\tfrac1{12}=\tfrac34\); then
\(\widetilde P_X=F(X)Z_X\le\tfrac34Z_X\), so
\(Z_X=(1-P_X)+\widetilde P_X\le1+\tfrac34Z_X\), i.e.\ \(Z_X\le4\).  With
\(\tfrac i{i-k}\le3\) this yields the estimate used below:
\begin{equation}\label{eq:line}
 \bigl|\widetilde P_X+h_kP_X-h_k\bigr|\le12\,\varepsilon
 \qquad\text{whenever } k=|X\cap[i-1]|\le\tfrac{2i}3 .
\end{equation}

\paragraph{Why two prefix counts are needed.}
A Boolean input constrains the two sums on its support only through the single
combination \(\widetilde P_X+h_kP_X\) in \eqref{eq:line}.  At one fixed
count, a change in \(P_X\) can therefore be offset by a change in
\(\widetilde P_X\).  In particular, the singleton Boolean input \(F(\{j\})\) does not
determine the baseline mass \(p_j\).  Using two counts changes the coefficient
\(h_k\) and separates the two terms.  Fix
\[
 k_1=\lfloor i/3\rfloor,
 \qquad
 k_2=\lfloor 2i/3\rfloor ;
\]
then \(3k_1\in\{i-2,i-1,i\}\) and \(3k_2\in\{2i-2,2i-1,2i\}\), and for
\(i\ge12\) clearing denominators gives
\begin{equation}\label{eq:dens}
 h_{k_1}\le\tfrac12,
 \qquad
 h_{k_2}\ge\tfrac32,
 \qquad
 \frac{i-1}{k_1}\le\frac{10}3,
 \qquad
 \frac{i-1}{k_2}\le\frac32,
\end{equation}
The inequalities are verified at the end of the proof.  In particular,
\(h_{k_2}-h_{k_1}\ge1\).  The swap step below also uses \(k_2\le i-2\),
which likewise holds for \(i\ge12\).  Throughout, \(J=[i-1]\) and
\(L=\{i+1,\dots,n\}\).

\paragraph{Comparing Boolean inputs with equal prefix counts.}
Two Boolean inputs with the same prefix count share the right side of
\eqref{eq:line}, so subtracting their lines bounds the difference of
the corresponding linear combinations by \(24\varepsilon\).  We first compare prefix
swaps and then compare inputs that differ on every future position.

Let \(j\ne j'\in J\), let \(k\in\{k_1,k_2\}\), and pick
any \(Y\subseteq J\setminus\{j,j'\}\) with \(|Y|=k-1\), possible since
\(k-1\le i-3\).  The Boolean inputs \(Y\cup\{j\}\) and \(Y\cup\{j'\}\) cancel
everything they share:
\[
 \bigl|(\tilde p_j-\tilde p_{j'})+h_k\,(p_j-p_{j'})\bigr|\le24\varepsilon .
\]
Subtracting the \(k_1\) and \(k_2\) versions (with possibly different
\(Y\) for each count) removes \(\tilde p_j-\tilde p_{j'}\), and dividing by
\(h_{k_2}-h_{k_1}\ge1\) leaves
\begin{equation}\label{eq:osc}
 \osc_{j\in J}\,p_j\le48\varepsilon .
\end{equation}

Next let \(Y\subseteq J\) with \(|Y|=k_2\).  The Boolean inputs \(Y\) and
\(Y\cup L\) also share their count, so
\(|\widetilde P_L+h_{k_2}P_L|\le24\varepsilon\).  Both terms are
nonnegative, and \(h_{k_2}\ge\tfrac32\), so
\begin{equation}\label{eq:fut}
 P_L\le16\varepsilon .
\end{equation}

\paragraph{Recovering the total prefix mass.}
Arrange \(J\) on a circle and let \(W_1,\dots,W_{i-1}\) be its cyclic
windows of length \(k\in\{k_1,k_2\}\); each \(j\in J\) lies in
\(k\) of them.  Summing \eqref{eq:line} over the windows and dividing by
\(k\) gives, with
\(h_k/k=1/(i-k)\),
\[
 \Bigl|\widetilde P_J+h_kP_J-\frac{i-1}{i-k}\Bigr|
 \le\frac{i-1}k\,12\varepsilon ,
\]
at most \(40\varepsilon\) for \(k_1\) and \(18\varepsilon\) for \(k_2\)
by \eqref{eq:dens}.  Subtracting the two displays cancels
\(\widetilde P_J\).  The remaining constant terms satisfy
\[
 \frac1{i-k_2}-\frac1{i-k_1}
 =\frac{k_2-k_1}{(i-k_1)(i-k_2)}
 =\frac{h_{k_2}-h_{k_1}}i ,
\]
so \(\bigl|(h_{k_2}-h_{k_1})\bigl(P_J-\frac{i-1}i\bigr)\bigr|
\le58\varepsilon\); dividing by \(h_{k_2}-h_{k_1}\ge1\),
\begin{equation}\label{eq:tot}
 \Bigl|P_J-\frac{i-1}i\Bigr|\le58\varepsilon .
\end{equation}

\paragraph{Completing the extraction.}
The mean of the prefix masses is \(P_J/(i-1)\), which by \eqref{eq:tot}
is within \(58\varepsilon/(i-1)\le6\varepsilon\) of \(1/i\), and each
\(p_j\) is within the oscillation \eqref{eq:osc} of that mean, so
\[
 \bigl|p_j-\tfrac1i\bigr|\le48\varepsilon+6\varepsilon=54\varepsilon
 \qquad(j\in J).
\]
Normalization gives \(p_i=1-P_J-P_L\), so \eqref{eq:tot} and
\eqref{eq:fut} imply
\(|p_i-\frac1i|\le58\varepsilon+16\varepsilon=74\varepsilon\).  For
\(j\in L\), \(0<p_j\le P_L\le16\varepsilon\), while \(T(i,j)=0\).
Altogether
\[
 \max_j\bigl|A(0)_{ij}-T(i,j)\bigr|\le C\varepsilon,
 \qquad C=74 .
\]

\paragraph{Verification of the numerical bounds.}
We have \(3k_1\in\{i-2,i-1,i\}\) and
\(3k_2\in\{2i-2,2i-1,2i\}\).  The first bound in \eqref{eq:dens} is
equivalent to \(3k_1\le i\).  For the second,
\[
 5k_2\ge\frac{5(2i-2)}3\ge3i
\]
when \(i\ge10\), which is equivalent to \(h_{k_2}\ge3/2\).  For the third,
\[
 10k_1\ge\frac{10(i-2)}3\ge3(i-1)
\]
when \(i\ge11\), giving the third bound.  The fourth follows from
\(3k_2\ge2i-2\).
Finally, \(k_2\le2i/3\le i-2\) when \(i\ge6\).  The constants in the
error bounds are \(12=3\cdot4\) in \eqref{eq:line},
\((\tfrac{10}3+\tfrac32)\cdot12=58\) in the window sums, and
\(C=58+16=74\) in the final step.
\end{proof}

\begin{proof}[Proof of Corollary~\ref{cor:adaptive}]
Set \(\varepsilon=n^{-p}\), \(N=n-1\), \(I=\{12,\ldots,n-1\}\), and
\(m=|I|=N-11\).

For the lower bound, fix any \(\varepsilon\)-accurate head with
\(d_h\le r\).  For all sufficiently large \(n\), every \(i\in I\) has
\(i\varepsilon\le n^{1-p}\le10^{-3}\), so Theorem~\ref{thm:adaptive} gives
\(\max_j|A(0)_{ij}-T(i,j)|\le74\varepsilon\); since
\(74\,i\varepsilon\le74\,n^{1-p}<1/2\) for large \(n\),
Lemma~\ref{lem:profile} applies to row \(i\) of \(A(0)\) at entrywise
error \(74\varepsilon\) and gives on \(I\) the uniform margins
\[
 \delta_i\le\log\frac{1+74\,i\varepsilon}{1-74\,i\varepsilon}
 \le222\,n^{1-p}=:\delta,
 \qquad
 \gamma_i\ge\log\frac{1-74\,i\varepsilon}{74\,i\varepsilon}
 \ge(p-1)\log n-\log222=:\gamma,
\]
using \(\log\frac{1+x}{1-x}\le3x\) and \(1-x\ge\frac13\) at
\(x=74\,i\varepsilon\le0.074\).
The all-zero score matrix \(S(0)\) satisfies \(\rho(S(0))\le r\) and
\(B(S(0))\le B_\infty\).  The lower-bound calculation in the proof of
Theorem~\ref{thm:sharp}, with \(N\) replaced by \(m\) and \(B\) by
\(B_\infty\), therefore gives
\[
 \log B_\infty\ge\frac mr
 \bigl((p-1)\log n+\log\log n+O_{p,r}(1)\bigr)+O_p(\log n).
\]
Replacing \(m=N-11\) by \(N\) changes only the remainder, and taking the
infimum over feasible heads proves the lower bound.

For the upper bound, take \(S_\star\) from Theorem~\ref{thm:upper}; its
statement gives \(\rank(S_\star)=\rho(S_\star)=r\).  Choose a rank
factorization \(S_\star=UV^\top\), let \(a_j,b_j\in\R^r\) be the \(j\)-th
rows of \(U,V\), and set
\[
 d_h=r,\quad d=2r,\quad e_j=(a_j,b_j),\quad u_j=0,
 \quad Q=[I_r\ \ 0],\quad K=[0\ \ I_r].
\]
Every payload realizes \(S(v)_{ij}=a_i^\top b_j=S_\star(i,j)\): the head
is content-independent and \(\varepsilon\)-accurate, with
\(B_\infty=B(S_\star)\).  Theorem~\ref{thm:upper}, evaluated at
\(\varepsilon=n^{-p}\) as in the proof of Theorem~\ref{thm:sharp}, gives
the matching upper estimate.
\end{proof}

The polynomial radius budget of Corollary~\ref{cor:budget} transfers to
the adaptive model.

\begin{corollary}[Adaptive radius budget]\label{cor:adaptivebudget}
Fix \(p>3/2\) and \(c>0\).  An \(\varepsilon\)-accurate head with
\(\varepsilon=n^{-p}\) and \(B_\infty\le n^c\) has
\[
 d_h\ge\left(\frac{p-3/2}{c+p-1}-o(1)\right)n .
\]
\end{corollary}

\begin{proof}
Set \(\varepsilon=n^{-p}\), \(N=n-1\), \(I=\{12,\ldots,n-1\}\), and
\(m=|I|=N-11\).  As in the proof of Corollary~\ref{cor:adaptive}, for all
sufficiently large \(n\), Theorem~\ref{thm:adaptive} and
Lemma~\ref{lem:profile} give on \(I\) the uniform margins
\[
 \delta=222\,n^{1-p},\qquad
 \gamma=(p-1)\log n-\log222,
\]
and the all-zero score matrix satisfies \(\rho(S(0))\le d_h\) and
\(B(S(0))\le B_\infty\le n^c\).

If \(d_h\ge m\), the conclusion is immediate because
\((p-3/2)/(c+p-1)<1\) and \(m/n\to1\).  Suppose instead that \(d_h<m\).
Eventually \(\gamma\ge\delta\), so Proposition~\ref{prop:hadamard}
gives \(2B_\infty\ge\gamma\ge\delta\) and
\[
 \gamma^m\le
 2^{m-1}(\delta\sqrt m)^{m-d_h}(2B_\infty\sqrt m)^{d_h},
\]
and therefore
\[
 d_h\log\frac{2B_\infty}{\delta}
 \ge m\log\frac{\gamma}{\delta\sqrt m}-(m-1)\log2.
\]
Uniformly in \(d_h\),
\[
 \log\frac{\gamma}{\delta\sqrt m}
 =(p-3/2)\log n+\log\log n+O_p(1),
 \qquad
 \log\frac{2B_\infty}{\delta}
 \le(c+p-1)\log n+O_{p,c}(1).
\]
The quantity
\(m\log(\gamma/(\delta\sqrt m))-(m-1)\log2\) is positive for large \(n\).
Replacing the multiplier of \(d_h\) by its upper bound, dividing by
\(m\log n\), and using \(m/n\to1\) proves the claim.  The remainders depend
only on \(p,c\), so the resulting \(o(1)\) is uniform in \(d_h\).
\end{proof}

\section{Experimental details}\label{app:experiment}

{\scriptsize\raggedright
The offline executable audit is
\path{experiments/run_one_head_experiments.py}; it writes
\path{experiments/output/one_head_bracket.csv},
\path{experiments/output/one_head_anchor.csv}, and
\path{experiments/output/one_head_thresholds.csv}, together with
\path{experiments/output/environment.txt}.\par}

The manuscript source, the three audit programs, their committed outputs,
and \path{C17/OneHead.lean} are available at
\url{https://anonymous.4open.science/r/rank-logit-law}.

The bracket grid uses
\(n\in\{12,16,20,24,32,40,48,56,64\}\) and
\(r\in\{1,2,4,8\}\).  The program constructs rank-one block factors from
binary64-generated budgets.  High-precision decimal arithmetic recomputes each
worst-row error and centered radius.  The program also solves
\eqref{eq:weighted} and evaluates \eqref{eq:upperbudget}.  Decimal precision follows the proved budget
with 80 guard digits.

The threshold audit generates the representability table and the head-dimension
crossings at \(n=8192\).  Its length searches record the last
admitted value followed by the first excluded value.  Dimension searches
record the last excluded value followed by the first admitted value.  The
environment file identifies the interpreter, platform, Decimal exponent
limit, and truncation cutoff used to produce the artifacts.  Each search
records its two-sided decision margin; the smallest reported is
\(2.3\times10^{-4}\) nats.  A second program,
\path{experiments/certify_one_head_thresholds.py} recomputes every reported
crossing and the \(n=8192\) anchor bound with integer Stirling numbers
and outward-rounded interval arithmetic, certifying each decision in
\path{experiments/output/one_head_interval_certificates.csv}.  The program
\path{experiments/verify_finite_inequalities.py} checks the permutation
count of Lemma~\ref{lem:count} exhaustively through \(m=8\) and inequalities
\eqref{eq:weighted}, \eqref{eq:constant}, and the Hadamard bound
directly on every bracket instance.

A Lean~4 development accompanies the project.  For the statements of this
paper it machine-checks the centered-rank-one, constant-profile
specialization of Theorem~\ref{thm:subset}: with the causal-difference
profile taken as a hypothesis, \(\gamma^N\le2B\,\delta^{N-1}\) for every
positive boundary count \(N\) (\path{C17/OneHead.lean}, building against
Mathlib alone).  The development builds with no \texttt{sorry} and no
project axiom, and its theorems depend only on \texttt{propext},
\texttt{Classical.choice}, and \texttt{Quot.sound}.  The general row-subset
inequality, the permutation count of Lemma~\ref{lem:count}, the Hadamard
bound, and the construction of Theorem~\ref{thm:upper} are not
formalized; their finite-instance verification is the executable audit
above.

\end{document}
